% ===== DRAFT: formalism (sec 0) and metrics (sec 4). Merged into answer.tex later. =====

\subsection*{1. Set-up: heterogeneous information networks, schemas and meta-paths}

\paragraph{Heterogeneous information network (HIN).}
Following Sun et al.~\cite{p2:pathsim} and the survey of Shi et al.~\cite{p2:hinsurvey}, a
\emph{heterogeneous information network} is a directed graph
\[
  G=(V,E,\phi,\psi),\qquad
  \phi:V\to\mathcal{A},\qquad
  \psi:E\to\mathcal{R},
\]
where every node $v\in V$ carries exactly one \emph{node type} $\phi(v)$ drawn from the type set
$\mathcal{A}$, every edge $e\in E$ carries exactly one \emph{relation type} $\psi(e)$ drawn from
$\mathcal{R}$, and
\[
  |\mathcal{A}|+|\mathcal{R}|>2 .
\]
The inequality is the definition of heterogeneity: if $|\mathcal{A}|=|\mathcal{R}|=1$ the network
collapses to the ordinary homogeneous graph treated in Problem~1. A relation
$R\in\mathcal{R}$ is typed at both ends, so it induces a map between two specific node types,
written $A_i\xrightarrow{R}A_j$; its inverse $R^{-1}$ is in general a different relation
($\textsf{writes}$ vs.\ $\textsf{written-by}$).

\paragraph{Network schema.}
The \emph{network schema} $T_G=(\mathcal{A},\mathcal{R})$ is the directed graph defined over the
node types $\mathcal{A}$ with edges the relation types $\mathcal{R}$. It is a meta-template: $G$ is
an instantiation of $T_G$ in the sense that every node and every edge of $G$ must map to some type
in $T_G$, so the schema tells you which kinds of connection are even expressible in the dataset.
This is exactly the object the question asks to be drawn for each benchmark, and it is the reason
the drawing carries information here and not in the homogeneous case: for Cora the schema is a
single node type with a single self-relation, whereas for DBLP it is a four-type star.

\paragraph{Meta-path.}
A \emph{meta-path} $\mathcal{P}$ of length $\ell$ is a path defined \emph{on the schema},
\[
  \mathcal{P}\;:\;A_1\xrightarrow{R_1}A_2\xrightarrow{R_2}\cdots\xrightarrow{R_\ell}A_{\ell+1},
\]
which denotes the composite relation $R=R_1\circ R_2\circ\cdots\circ R_\ell$ between type $A_1$ and
type $A_{\ell+1}$. A concrete node sequence $v_1v_2\cdots v_{\ell+1}$ in $G$ with
$\phi(v_i)=A_i$ and $\psi(v_iv_{i+1})=R_i$ is a \emph{path instance} of $\mathcal{P}$. A meta-path
is usually written by its type initials, so in DBLP $\mathrm{APA}$ abbreviates
$\text{Author}\xrightarrow{\textsf{writes}}\text{Paper}\xrightarrow{\textsf{written-by}}\text{Author}$.
Meta-paths matter because each one defines a different homogeneous projection of the same HIN, and
those projections carry different semantics: on DBLP, $\mathrm{APA}$ means \emph{co-authorship}
whereas $\mathrm{APCPA}$ means \emph{publishing at the same venue}, and two authors can be strongly
related under one and unrelated under the other. Essentially every method in
Section~\ref{p2:sec:baselines} is a different answer to the question ``how should the meta-path
structure be used?'' --- fixed and hand-chosen (metapath2vec, HAN), aggregated with attention
(MAGNN), learned (GTN), or dispensed with in favour of per-relation parameters (R-GCN, HGT,
simple-HGN).

% =====================================================================
\subsection*{4. Evaluation metrics}

Two conventions run through the whole heterogeneous literature and are worth stating before the
formulas. First, a HIN benchmark normally labels \textbf{exactly one node type} --- the author in
DBLP, the paper in ACM, the movie in IMDB --- so ``node classification accuracy on DBLP'' always
means accuracy on that one target type, and the other types exist only as context. Second, for
\emph{unsupervised} embedding methods the standard protocol is
\emph{embed $\to$ freeze $\to$ fit a linear classifier}: the embedding is learned without labels,
then held fixed while a logistic regression or linear SVM is trained on a fraction of the target
nodes, with the fraction varied over $20/40/60/80\%$ and mean $\pm$ std reported over several
splits. HAN, MAGNN and HeCo all report in this form.

\subsubsection*{4.1 Node classification}

Let $\mathcal{C}$ be the class set and, for class $c$, let $\mathrm{TP}_c,\mathrm{FP}_c,\mathrm{FN}_c$
be the true-positive, false-positive and false-negative counts. \textbf{Micro-F1} pools the
confusion matrix over all classes before computing the score:
\[
  \mathrm{Micro\text{-}P}=\frac{\sum_{c\in\mathcal{C}}\mathrm{TP}_c}
       {\sum_{c\in\mathcal{C}}\left(\mathrm{TP}_c+\mathrm{FP}_c\right)},\qquad
  \mathrm{Micro\text{-}R}=\frac{\sum_{c\in\mathcal{C}}\mathrm{TP}_c}
       {\sum_{c\in\mathcal{C}}\left(\mathrm{TP}_c+\mathrm{FN}_c\right)},
\]
\[
  \boxed{\;\mathrm{Micro\text{-}F1}
   =\frac{2\,\mathrm{Micro\text{-}P}\cdot \mathrm{Micro\text{-}R}}
          {\mathrm{Micro\text{-}P}+\mathrm{Micro\text{-}R}}\;}
\]
whereas \textbf{Macro-F1} computes a per-class F1 and averages them \emph{unweighted}:
\[
  \mathrm{F1}_c=\frac{2\,\mathrm{TP}_c}{2\,\mathrm{TP}_c+\mathrm{FP}_c+\mathrm{FN}_c},
  \qquad
  \boxed{\;\mathrm{Macro\text{-}F1}=\frac{1}{|\mathcal{C}|}\sum_{c\in\mathcal{C}}\mathrm{F1}_c\;}
\]
In the \emph{single-label} multi-class setting each node contributes exactly one false positive for
every false negative, so $\mathrm{Micro\text{-}P}=\mathrm{Micro\text{-}R}$ and Micro-F1 coincides
with plain accuracy --- which is why papers that report Micro-F1 rarely also report accuracy. The
pair is reported together because they answer different questions: Micro-F1 is dominated by the
frequent classes, Macro-F1 weights a rare class as heavily as a common one, and a large gap between
them is the signature of a method that has learned the class prior rather than the graph.

\paragraph{The multi-label caveat.} This equivalence \emph{fails} for the multi-label benchmarks. In
the HGB versions of IMDB and Freebase a node may carry several labels at once (a film is both
``Drama'' and ``Romance''), prediction is per-label thresholding rather than an $\arg\max$, and the
pooled confusion matrix is taken over the flattened node~$\times$~label indicator matrix. There
$\mathrm{Micro\text{-}P}\neq\mathrm{Micro\text{-}R}$ in general and Micro-F1 is \emph{not} accuracy.
Comparing a Micro-F1 from a multi-label IMDB against one from single-label DBLP is meaningless.

\paragraph{Other classification metrics.} \textbf{Accuracy} is reported mainly on OGB
(\texttt{ogbn-mag}'s official leaderboard metric is accuracy). \textbf{ROC-AUC} appears for binary
or strongly imbalanced targets. For the \emph{clustering} probe that HAN, MAGNN, DMGI and HeCo all
run --- $k$-means on the frozen embeddings with $k=|\mathcal{C}|$ --- the scores are normalised
mutual information and adjusted Rand index. With $Y$ the ground-truth partition and $\hat{C}$ the
predicted one,
\[
  \mathrm{NMI}(Y,\hat{C})=\frac{2\,I(Y;\hat{C})}{H(Y)+H(\hat{C})}\in[0,1],
  \qquad
  \mathrm{ARI}=\frac{\mathrm{RI}-\mathbb{E}[\mathrm{RI}]}{\max(\mathrm{RI})-\mathbb{E}[\mathrm{RI}]},
\]
where $I$ is mutual information, $H$ is entropy and $\mathrm{RI}$ is the Rand index; both are $0$
for a chance partition and $1$ for a perfect one, and ARI can go negative. These are the metrics
that test whether the embedding geometry is class-structured without a classifier being fitted at
all.

\subsubsection*{4.2 Link prediction}

The essential difference from the homogeneous case is that in a HIN link prediction is stated
\textbf{per relation type}: the task is not ``is there an edge'' but ``is there an edge
\emph{of type $R$}'' --- predict Author--Paper, or Paper--Venue, or User--Artist. Two consequences
follow. (i) \emph{Negatives must be sampled type-consistently}: a candidate negative for an
Author--Paper edge must join an author to a paper, never an author to a venue, because a
type-violating negative is trivially rejected by any type-aware model and inflates its score
against type-agnostic baselines. (ii) Scores are reported per relation and then averaged, so a
model can be strong on a dense relation and useless on a sparse one without that showing in a
single number.

Let the evaluation rank a positive edge against a set of sampled negatives. With
$\mathrm{rank}_i$ the position of the $i$-th true edge among its candidates and $Q$ the query set,
\[
  \boxed{\;\mathrm{MRR}=\frac{1}{|Q|}\sum_{i\in Q}\frac{1}{\mathrm{rank}_i}\;}
  \qquad
  \boxed{\;\mathrm{Hits}@K=\frac{1}{|Q|}\sum_{i\in Q}\mathbf{1}\!\left[\mathrm{rank}_i\le K\right]\;}
\]
so MRR is the mean reciprocal rank of the first true item and Hits@$K$ is the fraction of positive
edges ranked inside the top $K$ against their sampled negatives; $K\in\{1,3,10\}$ is conventional,
inherited from the knowledge-graph literature. \textbf{ROC-AUC} treats the problem as binary
scoring and equals the probability that a random positive outranks a random negative;
\textbf{average precision} (equivalently the area under the precision--recall curve) is
\[
  \mathrm{AP}=\sum_{n}\left(R_n-R_{n-1}\right)P_n ,
\]
with $P_n,R_n$ the precision and recall at the $n$-th threshold. AP is preferred to ROC-AUC when
the positive class is rare, which it always is for link prediction on a sparse graph --- ROC-AUC
is optimistic there because the huge true-negative count sits in its denominator.
\textbf{Precision@$K$} and \textbf{Recall@$K$} appear in the recommendation-flavoured HIN papers
(HERec on Yelp/Douban/MovieLens). HGB's link-prediction track standardises on \textbf{ROC-AUC and
MRR}.
